% Borrador de protocolo, versión 0.1, 22 de septiembre de 2026.
% No presenta como ejecutadas las pruebas comparativas pendientes.
% Requiere amsmath, amssymb y booktabs. Añadir referencias_capitulo6.bib
% a la bibliografía de la tesis; este archivo no incluye el preámbulo.
\chapter{Validación y experimentación}
\label{ch:validacion_experimentacion}

Este capítulo establece el procedimiento para comparar las arquitecturas
Qwen--CLIPS y GPT directo descritas en los capítulos anteriores. La evaluación
se centra en la transformación de instrucciones a objetivos
semánticos y planes de acciones del dominio GPSR. Se examinarán la precisión de
la interpretación, el cumplimiento de las tareas mediante el plan generado y
el rendimiento computacional de cada configuración.

El protocolo distingue el desarrollo del instrumento de evaluación de su
aplicación final. Las pruebas preliminares permitirán corregir los registros,
precisar las referencias y fijar la configuración experimental. Posteriormente,
la comparación se realizará sobre un conjunto reservado que no se utilizará
para ajustar los sistemas. Los resultados de esa comparación se presentarán en
el capítulo siguiente. 

\section{Diseño experimental}
\label{sec:diseno_experimental_cap6}

El diseño será pareado: cada orden del conjunto final se presentará a ambas
arquitecturas bajo condiciones iniciales y criterios de evaluación comunes.
La arquitectura completa será el factor principal de comparación. Cambian
simultáneamente el modelo, su especialización y el mecanismo de generación del
plan; por ello, las diferencias no se atribuirán exclusivamente al empleo de
reglas simbólicas.

La comparación responderá a cuatro preguntas: qué tan fielmente se recupera la
intención de la orden; qué proporción de planes satisface sus subtareas y
restricciones; cómo cambia ese comportamiento ante reformulaciones y mayor
composición; y cuál es el tiempo y consumo observable requerido para producir
la respuesta. La evaluación por capacidades se apoya en el enfoque de pruebas
de comportamiento de CheckList, que complementa las medidas agregadas con
casos dirigidos a comportamientos lingüísticos específicos
\cite{ribeiro2020checklist}. La distinción entre una respuesta plausible y un
plan verificable se relaciona con la evaluación de acciones y cambios de
estado abordada por PlanBench \cite{valmeekam2023planbench}. En esta tesis se
desarrollará una rúbrica propia para el contrato operativo de Justina.

\subsection{Conjuntos de desarrollo y prueba}
\label{subsec:conjuntos_cap6}

Se distinguirán tres usos de los datos. El corpus de 40\,000 órdenes empleado
en la especialización de Qwen, dividido en entrenamiento y validación, conserva
la función descrita en el capítulo de implementación. El conjunto de
desarrollo permitirá preparar el experimento. El conjunto de prueba final se
reservará para la comparación con las configuraciones ya fijadas. Las
comprobaciones funcionales previas y las predicciones efectuadas sobre la
partición de validación no sustituirán este último conjunto.

El conjunto de desarrollo preparado contiene 60 órdenes, distribuidas como
se muestra en la Tabla~\ref{tab:desarrollo_cap6}. Las tres formulaciones de una
misma intención se mantendrán agrupadas mediante un identificador común.

\begin{table}[htbp]
\centering
\small
\caption{Organización del conjunto de desarrollo propuesto.}
\label{tab:desarrollo_cap6}
\begin{tabular}{|p{0.56\textwidth}|r|r|}
\toprule
Bloque & Órdenes & Intenciones \\
\midrule
Funcional y lingüístico: formulación canónica, paráfrasis y ruido leve & 48 & 16 \\
Límites de representación y composición & 6 & 6 \\
Ambigüedad, contradicción y tareas fuera del repertorio & 6 & 6 \\
\midrule
Total & 60 & 28 \\
\bottomrule
\end{tabular}
\end{table}

Estos 60 casos no constituyen 60 intenciones independientes ni acreditan
cobertura exhaustiva de las familias del generador. Cincuenta y cuatro fichas
contienen objetivos propuestos y seis describen conductas de aclaración o
rechazo. La existencia de objetivos no implica que todos los casos pertenezcan
al bloque común resoluble: las capacidades cuyo contrato sea incompleto,
como determinadas formas de seguimiento o identificación, requieren revisión.

Cada ficha incluirá la orden, los objetivos de referencia
\(G_{\mathrm{ref}}\), las subtareas requeridas, las restricciones, los
comportamientos inadmisibles y las equivalencias aceptables. Las referencias
se revisarán a partir de la intención expresada y del contrato de acciones,
antes de observar las salidas de los modelos.

Se auditarán las coincidencias con el corpus completo mediante igualdad
literal, limpieza básica y aplicación del normalizador utilizado por Qwen.
Una coincidencia con las 40\,000 filas no identifica por sí misma
la partición a la que perteneció la fila. Asimismo, ausencia de coincidencias
bajo estas transformaciones no demuestra novedad semántica, independencia de
plantilla ni ausencia en el preentrenamiento de los modelos.

En desarrollo podrán conservarse coincidencias identificadas como casos
conocidos, porque su función es diagnosticar el procedimiento. No se modificarán
indefinidamente las órdenes sólo para obtener un contador de coincidencias
igual a cero. Para la prueba final se excluirán los duplicados detectados y
las instancias de intención utilizadas durante el desarrollo, junto con sus
reformulaciones. Se podrán conservar las mismas clases de tarea para evaluar
generalización dentro del dominio; esto no se describirá como evaluación de
familias nunca vistas.

El tamaño definitivo del conjunto final se establecerá después del piloto y
antes de ejecutar la comparación. Como propuesta de cobertura, se consideran
diez intenciones por familia admisible y tres formulaciones por intención.
Si \(F\) es el número de familias incluidas, se obtendrían \(10F\) intenciones
y \(30F\) órdenes. Incluir las 29 familias del generador produciría 290
intenciones y 870 órdenes, pero esta cifra queda condicionada a la revisión de
su correspondencia con las capacidades y al presupuesto de ejecución. No
representa un tamaño calculado mediante potencia estadística.

La admisión al bloque principal dependerá de que la tarea esté suficientemente
especificada y pueda evaluarse mediante el contrato declarado, no de que ambos
sistemas consigan resolverla. Una arquitectura que falle en una tarea
admisible conservará ese fallo. Las tareas cuyo significado o verificación no
pueda establecerse se identificarán antes de la comparación y se estudiarán
como diagnósticos. Los casos ambiguos y fuera de dominio formarán un bloque
separado, con cantidad y distribución fijadas también antes de las ejecuciones.

\subsection{Procedimiento experimental}
\label{subsec:procedimiento_cap6}

El procedimiento se organizará en preparación, piloto, fijación de la
configuración y evaluación final. Durante la preparación se revisarán las
fichas y se construirá un registro común para ambas arquitecturas. Los
validadores de formato existentes se reutilizarán donde corresponda, pero se
añadirán verificaciones de contenido basadas en las subtareas y restricciones
de cada caso. Los registros actuales de consola y planes no constituyen por sí
solos este instrumento completo.

El evaluador se comprobará con ejemplos correctos y con errores introducidos
deliberadamente: omitir una adquisición, cambiar el destinatario, eliminar un
atributo o truncar la secuencia. Esta comprobación permitirá detectar un
evaluador que acepte cualquier plan bien formado o que rechace alternativas
válidas por no reproducir una única secuencia textual. Las reglas de
puntuación se fijarán antes de observar el conjunto final.

El piloto comenzará con diez casos de desarrollo que cubran operaciones
básicas y referencias; después se completarán los 60 casos. Los primeros diez
forman parte de esos 60 y no constituyen datos adicionales. El objetivo será
confirmar trazabilidad, reinicio de estado, límites de salida, tiempos de
espera y aplicación de la rúbrica. Las repeticiones del piloto servirán para
estimar variabilidad y coste, sin presentarse como resultados finales.

En cada ejecución se seguirá esta secuencia:
\begin{enumerate}
\item Recuperar la orden y su identificador del manifiesto experimental. La
arquitectura recibirá únicamente el texto de la orden, además de su
configuración fija; las referencias permanecerán en el evaluador.
\item Restablecer el contexto del ensayo. Se utilizará una solicitud sin
historial conversacional y, para CLIPS, un entorno sin hechos de la orden
anterior. La configuración inicial prevista será ubicación desconocida,
ningún objeto sostenido, foco desactivado y manipulación habilitada.
\item Iniciar la medición del recorrido completo y procesar la orden mediante
la arquitectura correspondiente. Se conservarán las representaciones
intermedias disponibles, sin sustituir la salida original por una corrección
manual.
\item Esperar una respuesta terminada o registrar un fallo o agotamiento del
tiempo establecido.  
\item Conservar la respuesta cruda, comprobar el contrato de salida y detener
la medición de latencia de respuesta. La evaluación semántica posterior y la
revisión humana se realizarán fuera de ese intervalo.
\item Evaluar objetivos, plan y dependencias con la misma rúbrica, registrar
las incidencias y asociar el resultado con la versión de la orden y de la
configuración.
\end{enumerate}

Se evitará mezclar respuestas de órdenes diferentes. Como la publicación de
planes actual no proporciona por sí sola un identificador de solicitud, el
instrumento deberá incorporar correlación explícita o aislamiento verificable
por ensayo. Después de un tiempo agotado se cancelará o aislará el trabajo
pendiente antes de enviar otra orden. Vaciar una variable local sin controlar
publicaciones tardías no basta para garantizar la correspondencia.

El manifiesto conservará las huellas de las entradas, referencias, adaptador,
prompt, esquemas, normalizador, catálogos y reglas; las revisiones del código;
las versiones del entorno; los parámetros de generación; los tiempos de
espera; y la política de reintentos. Tras el piloto se fijarán estos elementos.
Si posteriormente se corrige un componente usando errores del conjunto final,
la nueva configuración se distinguirá como una evaluación posterior y no
sustituirá silenciosamente la comparación original.

\subsection{Repeticiones y control de variabilidad}

Se propone realizar tres repeticiones independientes por orden y arquitectura
en la evaluación final. El número se confirmará mediante la estimación de
coste y duración del piloto, antes de ejecutar el conjunto reservado. No se
escogerá la mejor respuesta de las repeticiones. Qwen conservará generación
sin muestreo; ello no elimina la necesidad de registrar variación de latencia
y posibles diferencias de la infraestructura.

El orden de las condiciones se aleatorizará por bloques con una semilla
registrada, alternando qué arquitectura procesa primero cada orden. Las
ejecuciones serán secuenciales para evitar interferencia entre solicitudes y
atribución ambigua de planes. Los identificadores de repetición vincularán
los resultados pareados, aunque no impliquen compartir una semilla de
generación entre proveedores.

Los reintentos no serán correcciones selectivas de respuestas erróneas. Se
propone conservar la primera tentativa como observación principal y registrar
cualquier reintento de infraestructura por separado. Un error de servicio o
tiempo agotado se mantendrá en el denominador de disponibilidad y de éxito de
extremo a extremo. El análisis condicionado a respuestas recibidas se
identificará como secundario.

Las variantes y repeticiones de una misma intención están relacionadas. Para
comparar tasas se promediarán primero los resultados de cada intención y se
calculará la diferencia pareada entre arquitecturas. Se propone estimar un
intervalo de confianza del 95\% mediante remuestreo de intenciones completas,
manteniendo dentro de cada grupo sus variantes, repeticiones y ambos sistemas.
El remuestreo se estratificará por familia cuando el diseño final lo requiera;
se registrarán su semilla y número de remuestras. No se tratarán todas las
variantes como observaciones independientes.

\section{Condiciones experimentales}
\label{sec:condiciones_cap6}

Ambas arquitecturas recibirán las mismas órdenes originales en inglés y se
evaluarán con el mismo vocabulario semántico, catálogos, repertorio de acciones
y convenciones de cierre. El normalizador de Qwen permanecerá como parte de
su arquitectura; GPT conservará el procesamiento establecido por su prompt.
La comparación principal medirá esas configuraciones completas. Una eventual
ablación del normalizador se definirá como análisis adicional, sin mezclar sus
resultados con la condición principal.

\subsection{Arquitectura Qwen--CLIPS}

La condición Qwen--CLIPS utilizará el adaptador especializado de
Qwen2.5-7B seleccionado mediante pérdida de validación. El entrenamiento
registrado seleccionó el checkpoint 5000 y guardó el adaptador restaurado en
\texttt{nl2cd\_qwen7b}. Para cada campaña se verificará la identidad del
artefacto cargado por el servidor, de modo que una modificación posterior de
esa carpeta no altere inadvertidamente la condición.

La cadena experimental será:
\begin{equation}
x \longrightarrow G_{\mathrm{pred}}^{\mathrm{Qwen}}
\longrightarrow \operatorname{Adaptar}
       (G_{\mathrm{pred}}^{\mathrm{Qwen}})
\longrightarrow P^{\text{Qwen--CLIPS}}.
\end{equation}
Se conservarán el texto normalizado, los objetivos predichos antes de su
adaptación, las transformaciones realizadas, los hechos insertados y el plan.
Los objetivos originales se utilizarán para evaluar la interpretación del
modelo; los hechos permitirán examinar pérdidas o cambios introducidos en la
integración.

El límite actual de generación y los tiempos de espera se comprobarán con las
órdenes compuestas del piloto. Si se ajustan, se documentarán los nuevos
valores antes de la prueba final. No se impondrá igualdad numérica de tokens
con GPT, porque Qwen produce objetivos y GPT produce además el plan; se
establecerán límites adecuados a cada contrato y se informará el truncamiento.

\subsection{Arquitectura GPT directo}

La condición GPT utilizará el programa independiente
\texttt{planner\_test\_node.py}, el prompt del dominio y el esquema de respuesta
estructurada. El modelo configurado actualmente es \texttt{gpt-5.6-terra}; se
registrarán el identificador solicitado, el informado por el servicio y la
fecha de ejecución. El recorrido será:
\begin{equation}
x \longrightarrow
\bigl(G_{\mathrm{pred}}^{\mathrm{GPT}},P^{\mathrm{GPT}}\bigr).
\end{equation}
Los objetivos y el plan se producen en una misma respuesta. No se enviarán
los objetivos de GPT a CLIPS para construir o corregir su plan.

La independencia del programa es compatible con la comparación, porque el
producto evaluado es un plan representado bajo el contrato de Justina. No es
necesario conectar GPT al ejecutor físico para medir ese producto. Sí será
necesario aplicar las mismas condiciones iniciales declaradas y la misma
rúbrica de contenido. La latencia incluirá el acceso al servicio externo y
se interpretará como rendimiento del sistema desplegado.

El prompt comunica convenciones operativas próximas a las utilizadas por
CLIPS. Por tanto, la comparación corresponde a generación directa guiada por
esas instrucciones frente a expansión mediante reglas. No se interpretará
como una prueba de planificación sin conocimiento del dominio ni como un
experimento que aísla únicamente el efecto de CLIPS.

\subsection{Control con objetivos de referencia y CLIPS}

La condición de control reemplazará la predicción de Qwen por
\(G_{\mathrm{ref}}\) y conservará el mismo adaptador y las mismas reglas de
la cadena principal:
\begin{equation}
G_{\mathrm{ref}} \longrightarrow
\operatorname{Adaptar}(G_{\mathrm{ref}})
\longrightarrow P^{\mathrm{control}}.
\end{equation}
Se utilizará el adaptador del recorrido de ejecución, no otro conversor que
produzca hechos diferentes. Se reiniciará el estado bajo las mismas
condiciones. Este control no constituye una tercera arquitectura competidora
y su latencia no se comparará como si incluyera interpretación lingüística.

Si el control satisface una orden y Qwen--CLIPS no, la inspección de los
objetivos y hechos ayudará a localizar la diferencia. Si el control también
falla, se revisarán la referencia, la adaptación, las reglas y el contrato de
acciones. La coincidencia de fallos no prueba por sí sola una causa concreta.
El plan de control se evaluará con la rúbrica independiente; no se convertirá
automáticamente en la respuesta correcta contra la cual juzgar a GPT.

\section{Evaluación de la interpretación semántica}
\label{sec:semantica_cap6}

\subsection{Exactitud de objetivos}

Se distinguirán coincidencia exacta y equivalencia semántica. La primera
requerirá que la lista de objetivos predicha coincida con la referencia
canónica, conservando orden y parámetros. Para el conjunto resoluble,
\begin{equation}
\mathrm{EM}_{G}^{a}=
\frac{1}{NR}\sum_{i=1}^{N}\sum_{r=1}^{R}
\mathbf{1}\!\left[G_{\mathrm{pred},i,r}^{a}=G_{\mathrm{ref},i}\right],
\end{equation}
donde \(a\) identifica la arquitectura, \(N\) el número de órdenes y \(R\)
el número fijado de repeticiones. Una salida ausente o no interpretable no
constituirá coincidencia. Las órdenes ambiguas o fuera de dominio se
evaluarán separadamente y no tendrán una lista vacía como referencia correcta.

La equivalencia semántica permitirá variantes previamente justificadas que
conserven subtareas, entidades y relaciones. Por ejemplo, un sinónimo de una
propiedad de tamaño puede diferir textualmente sin cambiar la intención. No
se permitirá eliminar atributos, intercambiar destinatarios ni reordenar
dependencias para aumentar artificialmente la coincidencia. Las decisiones
de equivalencia se documentarán con reglas comunes para ambos sistemas.

\subsection{Validez estructural}

Se verificará que la salida pueda analizarse, que los objetivos utilicen el
vocabulario final y que sus argumentos cumplan el esquema. Los doce objetivos
considerados son \texttt{go}, \texttt{find}, \texttt{take}, \texttt{deliver},
\texttt{place}, \texttt{guide}, \texttt{follow}, \texttt{count}, \texttt{tell},
\texttt{save}, \texttt{answer\_question} y \texttt{greet}. Un esquema aceptado
no garantiza correspondencia con la orden. Se conservarán separadas la
validez de estructura, la corrección semántica y la disponibilidad de respuesta.

\subsection{Entidades y parámetros}

La revisión comprobará objeto, persona, ubicación de origen, destino,
destinatario, tipo de entidad y atributos exigidos. Se distinguirán omisiones,
sustituciones y entidades inventadas. La medida agregada de orden completamente
correcta se acompañará de tasas por campo cuando haya suficientes casos.

Las secuencias de distinta longitud requieren alineación de subtareas: no se
comparará ciegamente el objetivo de una posición con el de la misma posición
en una lista desplazada. Las omisiones y adiciones conservarán su penalización
y la alineación no podrá ocultar cambios de orden relevantes. Las tasas de
acierto sobre campos evaluables informarán su denominador y se acompañarán
del resultado sobre todas las órdenes.

\subsection{Dependencias entre objetivos}

Se verificará la conservación de referencias y precedencias: buscar antes de
adquirir, adquirir antes de entregar, mantener el mismo objeto durante el
transporte y asociar la interacción con la persona solicitada. Las fichas
describirán estas relaciones sin exigir que todo orden alternativo sea
idéntico al de la referencia si respeta las mismas dependencias.

Un ejemplo de desarrollo es \emph{Bring me a mug from the cabinet}. La
referencia propuesta contiene:
\begin{verbatim}
go(cabinet)
find(mug, kind=object)
take(mug)
deliver(mug, to=me)
\end{verbatim}
Un resultado que cambie \texttt{cabinet} por otra ubicación tendrá un error de
interpretación aunque produzca después un plan bien formado. Si los objetivos
son correctos pero el plan omite la adquisición, el fallo aparecerá en la
etapa de planificación. Este ejemplo es ilustrativo y no describe una salida
experimental observada.

\section{Evaluación de la planificación}
\label{sec:planificacion_cap6}

La evaluación del plan utilizará subtareas y restricciones, no igualdad
literal con una única lista de acciones. Se admitirán anuncios equivalentes
y referencias internas consistentes, preservando los efectos exigidos. La
revisión automatizada se complementará con inspección humana cuando la
equivalencia no quede resuelta por los verificadores. Se ocultará la etiqueta
de arquitectura durante esa revisión en la medida en que el formato lo permita.

\subsection{Acciones válidas}

Se comprobarán nombres de acciones, número y tipo de argumentos, entidades
admisibles, numeración y convenciones de inicio y cierre. La separación entre
anuncios y operaciones deberá respetar el contrato compartido. Una respuesta
con JSON válido o con un marcador de finalización no se considerará un plan
correcto únicamente por esas propiedades.

\subsection{Consistencia del plan}

Se recorrerá la secuencia mediante un estado simbólico mínimo que represente
ubicación prevista, posesión del objeto, referente de interacción y datos
adquiridos cuando corresponda. El contrato del ejecutor determinará qué
acciones modifican esas variables. Este recorrido es un verificador de
consistencia, no una simulación de percepción o dinámica física.

Los mensajes hablados ordinarios no acreditarán por sí solos adquirir un
objeto, seguir a una persona o comprobar una identidad. Si una clave de
protocolo activa una operación en el ejecutor, será necesario documentar ese
efecto antes de concederle cumplimiento. Cuando no pueda precisarse, el
resultado se marcará como no verificable con el contrato disponible.

\subsection{Cumplimiento de subtareas}

Para cada orden se comprobará si el plan cubre todas las subtareas y respeta
las restricciones obligatorias. La variable principal será el cumplimiento
completo del plan en el bloque principal previamente definido:
\begin{equation}
C_{i,r}^{a}=
\begin{cases}
1, & \text{si el plan es válido y satisface todas las exigencias;}\\
0, & \text{si hay incumplimiento o no se obtiene un plan evaluable.}
\end{cases}
\end{equation}
La tasa correspondiente será
\begin{equation}
\mathrm{Cumplimiento}^{a}=
\frac{1}{NR}\sum_{i=1}^{N}\sum_{r=1}^{R}C_{i,r}^{a}.
\end{equation}
Se informarán también las fracciones de subtareas cumplidas para localizar
fallos parciales, sin sustituir la medida de orden completa. Un indicador
secundario exigirá simultáneamente interpretación correcta y plan correcto,
de modo que un plan acertado acompañado de objetivos incorrectos no oculte
una inconsistencia entre productos.

La categoría \emph{no verificable} se conservará separada de \emph{incumple}
en el diagnóstico. Si surge en un caso del bloque principal después de fijar
el protocolo, no se eliminará silenciosamente del denominador: se informará
como éxito no acreditado, se explicará la carencia del evaluador y se
presentará un análisis de sensibilidad. La proporción de casos verificables
permitirá valorar el alcance efectivo de la medida.

\subsection{Precondiciones y dependencias}

El verificador comprobará que los efectos requeridos estén disponibles antes
de cada acción dependiente. Para la orden de la taza, no bastará con enumerar
una entrega: el plan deberá localizar y adquirir la taza, conservar su
referencia y dirigir la entrega al operador. Cuando el contrato de entrega
incluya movimientos de manera implícita, se reconocerán esos efectos
documentados sin exigir una acción redundante de navegación.

Las condiciones iniciales serán supuestos declarados del ensayo. La
existencia y detectabilidad física de los objetos o personas no se comprobará
en este experimento. Tampoco se asumirán efectos simplemente porque CLIPS
actualice un hecho interno: el estado previsto del plan se contrastará con
el significado del contrato de acciones.

\subsection{Planificabilidad y terminación}

Se registrará si la cadena produce un plan terminado bajo su contrato y si
quedan objetivos pendientes o diagnósticos de falta de soporte. En CLIPS se
examinará la terminación efectiva del procesamiento, no sólo la recepción de
un mensaje no vacío. Una regla alternativa que consuma un objetivo sin
satisfacerlo podrá permitir terminar, pero no acreditará cumplimiento.

Se diferenciarán cuatro hechos: respuesta recibida, formato válido,
terminación del proceso y cumplimiento de la orden. Esta separación permitirá
explicar planes que concluyen formalmente y omiten tareas, así como fallos de
infraestructura que impiden observar el resultado del planificador.

\section{Clasificación de errores}
\label{sec:errores_cap6}

\subsection{Errores de interpretación}

Se clasificarán como errores de interpretación las subtareas omitidas o
inventadas, las entidades incorrectas, la pérdida de atributos, la resolución
errónea de pronombres y el cambio de dependencias en los objetivos predichos.
En Qwen se examinará también si el normalizador altera la orden antes de la
inferencia. Una modificación errónea del texto se distinguirá del fallo de
generación del modelo dentro de la cadena completa.

\subsection{Errores de adaptación y planificación}

Los errores de adaptación corresponderán a diferencias entre los objetivos
producidos y los hechos insertados, por ejemplo pérdida de un atributo o
reasignación de un tipo de entidad. Los errores de planificación incluirán
acciones inválidas, precondiciones incumplidas, referencias inconsistentes y
subtareas no materializadas pese a contar con objetivos suficientes.

También se registrarán separadamente truncamiento, ausencia de respuesta,
agotamiento del tiempo, fallos de servicio y problemas de asociación entre
solicitud y plan. No se atribuirá automáticamente a razonamiento un fallo
cuyo origen sólo sea observable como error técnico.

\subsection{Errores combinados y alcance de la atribución}

Una ejecución podrá recibir más de una etiqueta. Por ejemplo, una entidad
incorrecta en los objetivos y una entrega sin adquisición son errores
distintos. Se registrará el primer punto de divergencia observable y las
incidencias posteriores, sin asumir que todas comparten una causa.

La condición de control permitirá examinar la cadena simbólica con referencias
revisadas. En GPT, objetivos y plan se generan conjuntamente; su comparación
permite detectar inconsistencias entre productos, pero no demuestra una
separación causal interna entre interpretación y planificación. Las tablas de
errores conservarán esta diferencia de observabilidad entre arquitecturas.

\section{Robustez, límites del dominio y alcance físico}
\label{sec:robustez_cap6}

\subsection{Variación lingüística y composición}

Las variantes canónica, parafraseada y con ruido leve conservarán una misma
intención revisada. Se comparará el cumplimiento entre variantes y se
registrará cuántos grupos mantienen éxito en todas ellas. Obtener tres
respuestas equivalentes pero incorrectas no se interpretará como robustez
satisfactoria; la estabilidad se informará junto con la corrección.

Los casos compuestos permitirán evaluar conservación de referentes y
dependencias al aumentar las subtareas. El análisis se organizará por número
de objetivos, entidades y cambios de referente. Un resultado peor en órdenes
largas no se atribuirá únicamente a su longitud cuando también cambien familia
o dificultad. Se mantendrá una descripción de esas diferencias.

\subsection{Ambigüedad, restricciones y capacidades no verificables}

Se examinarán órdenes con referentes ausentes, destinos no especificados,
acciones fuera del repertorio y contradicciones. Para GPT se evaluará si la
respuesta solicita la información necesaria o explica adecuadamente el
rechazo, sin generar un plan basado en supuestos inventados. Se separará un
rechazo del dominio de un rechazo de la API o error del servicio.

Qwen no implementa el mismo contrato conversacional de aclaración y rechazo.
Se registrará su conducta observable ---bloqueo, error, objetivos propuestos o
plan generado--- sin convertir un error de formato en una aclaración
correcta. Estos resultados describirán cobertura funcional y límites; no se
mezclarán con la tasa principal de tareas resolubles.

Los diagnósticos incluirán atributos conjuntos de personas, conteos filtrados,
seguimiento, identificación por nombre y propiedades que no puedan
determinarse con los efectos disponibles. Por ejemplo, tamaño y peso no se
considerarán intercambiables. Un plan que use una aproximación no autorizada
se distinguirá de uno que satisface la propiedad solicitada. Las referencias
no se modificarán para aceptar una limitación común de las dos arquitecturas.

\subsection{Delimitación de la evaluación física}

El protocolo principal no requiere accionar al robot. Se detendrá el recorrido
en la obtención y verificación del plan; los nodos ejecutores físicos no
participarán en la campaña. Esta delimitación permitirá completar la
comparación aun cuando la plataforma presente problemas técnicos.

Las conclusiones se referirán a fidelidad semántica, consistencia simbólica y
rendimiento computacional. No permitirán estimar probabilidad de completar
una tarea en un entorno real ni certificar ausencia de daños. Si más adelante
se realizan demostraciones físicas, se describirán como una etapa adicional
con condiciones, casos y métricas propios. No se exigirá simetría física entre
una arquitectura integrada con ROS y un programa de evaluación independiente
para responder las preguntas de este protocolo.

\section{Rendimiento computacional}
\label{sec:rendimiento_cap6}

\subsection{Latencia}

La medida principal de latencia comenzará antes de enviar la orden original
y terminará cuando el evaluador reciba una respuesta terminal y complete la
validación automática del contrato. Para Qwen--CLIPS incluirá comunicación,
normalización, inferencia, adaptación y expansión; para GPT incluirá la
solicitud remota, generación, recepción y validación. La calificación semántica
detallada posterior quedará fuera de este intervalo.

Se mantendrá cargado el modelo local para medir operación habitual. La carga
inicial, la puesta en marcha de servicios y las solicitudes de calentamiento
se registrarán por separado y utilizarán casos de desarrollo, no entradas de
prueba final. Se fijará una única modalidad de despliegue de Qwen para la
comparación principal; una modalidad remota adicional se identificará como
otra condición de infraestructura.

Se presentarán mediana y percentil 95, junto con el número de observaciones,
la tasa de respuestas recibidas y la de planes correctos. Los tiempos hasta
fallo se informarán aparte; un fallo rápido no se interpretará como una
ventaja de latencia. Los tiempos agotados se conservarán como observaciones
censuradas al umbral declarado, sin asignarles un tiempo de éxito inventado.
Los tiempos internos de Qwen o CLIPS servirán para diagnóstico, pero no se
compararán directamente con la latencia total del servicio GPT.

\subsection{Uso de recursos}

En la plataforma local se registrarán memoria RAM y memoria GPU observables,
especificando si la medida corresponde al proceso, a memoria asignada por el
framework o al dispositivo completo. Se indicarán la frecuencia de muestreo,
la carga residente y el máximo observado durante la campaña. Un máximo
acumulado desde el inicio del proceso no se presentará como un máximo exclusivo
de una solicitud.

Para GPT no se dispone de acceso a los recursos internos del proveedor. La
memoria del cliente local no representa la memoria necesaria para ejecutar
el modelo remoto. Por tanto, se presentarán por separado los requisitos
locales de Qwen--CLIPS y los recursos observables del acceso al servicio GPT,
sin afirmar una comparación directa de consumo interno o energía.

\subsection{Tokens y coste del modelo GPT}

Se conservará el uso de tokens informado por el servicio para cada solicitud,
incluyendo el desglose disponible de entrada, salida y caché. Los tokens de
razonamiento se interpretarán según su pertenencia a los totales reportados,
sin sumarlos dos veces. El coste se obtendrá con la tarifa o facturación
aplicable en la fecha de ejecución, cuya fuente y moneda quedarán registradas.
Si no existe una tarifa verificable para el acceso utilizado, se informarán
tokens y cargo observado, sin inventar un coste teórico.

El informe distinguirá coste total de la campaña y coste por orden con plan
correcto. Los reintentos y solicitudes fallidas con consumo facturado se
conservarán en el total. Para Qwen, ausencia de un cargo por solicitud no
equivale a coste computacional nulo; no se construirá una comparación monetaria
integral sin medir y justificar los componentes adicionales.

La fijación de datos, referencias, configuraciones y denominadores permitirá
relacionar los resultados del capítulo siguiente con un procedimiento
explícito. El número final de casos, repeticiones y parámetros operativos se
incorporará cuando concluya el piloto, conservando los registros que
justifiquen esas decisiones.
